%------------------------------------------------------------------------------%
\documentclass[fleqn,utf8,aspectratio=169,12pt]{beamer}

\usepackage{gaal}

\title{Subespaços Vetoriais}

%------------------------------------------------------------------------------%
\begin{document}

\addtitle

%------------------------------------------------------------------------------%
\section{Definição}

%------------------------------------------------------------------------------%
\begin{frame}{Subespaço}
  Dado um espaço vetorial $V$

  \pause
  \medskip
  Um \emph{subconjunto} $W\subseteq V$
  \pause
  é um \emph{subespaço} de $V$

  \pause
  \medskip
  quando $W$ também é um \emph{espaço vetorial} com as operações de $V$
\end{frame}

%------------------------------------------------------------------------------%
\begin{frame}{Critério de Subespaço}
  Um subconjunto $W\subseteq V$ é um subespaço se
  \[
    \pause
    \vec{0}\in W
  \]

  \[
    \pause
    \vec{u}, \vec{v} \in W
    \pause
    \quad\Rightarrow\quad
    \vec{u} + \vec{v} \in W
  \]

  \[
    \pause
    \vec{u} \in W
    \pause
    \quad a \in \R
    \pause
    \quad\Rightarrow\quad
    a\vec{u} \in W
  \]
\end{frame}

%------------------------------------------------------------------------------%
\begin{frame}{Critério Alternativo}
  Basta verificar
  \[
    \pause
    \vec{0} \in W
  \]
  \pause
  \[
    a\vec{u}+b\vec{v}\in W
    \pause
    \qquad
    \text{ para quaisquer }
    \vec{u},
    \vec{v}\in W
    \quad
    a, b\in\R
  \]
\end{frame}

%------------------------------------------------------------------------------%
\begin{frame}{Subespaços Triviais}
  Todo espaço vetorial possui dois subespaços
  \[
    \pause
    W_1 = \{\vec{0}\}
  \]
  \pause e
  \[
    W_2 = V
  \]
\end{frame}

%------------------------------------------------------------------------------%
\section{Subespaços de $\R^n$}

%------------------------------------------------------------------------------%
\begin{frame}{Subespaços de $\R^2$}
  Os subespaços de $\R^{2}$ podem ser
  \[
    \pause
    W_1 = \{\vec{0}\}
  \]

  \[
    \pause
    W_2 = \text{uma reta que passa pela origem}
  \]

  \[
    \pause
    W_3 = \R^2
  \]
\end{frame}

%------------------------------------------------------------------------------%
\begin{frame}{Subespaços de $\R^{3}$}
  Os subespaços de $\R^{3}$ podem ser
  \[
    \pause
    W_1 = \{\vec{0}\}
  \]

  \[
    \pause
    W_2 = \text{uma reta que passa pela origem}
  \]

  \[
    \pause
    W_3 = \text{um plano que passa pela origem}
  \]

  \[
    \pause
    W_4 = \R^{3}
  \]
\end{frame}

%------------------------------------------------------------------------------%
\section{Subespaços Gerados}

%------------------------------------------------------------------------------%
\begin{frame}{Subespaço Gerado}
  Dado um conjunto de vetores de um espaço vetorial $V$
  \[
    G = \left\{ \vec{v}_{1}, \ldots, \vec{v}_{k} \right\}
  \]

  \bigskip
  \pause
  O conunto dos vetores gerados por $G$
  \[
    \pause
    W = \ger(G) = \ger \left\{ \vec{v}_{1}, \ldots, \vec{v}_{k} \right\}
  \]
  \pause
  é um subespaço de $V$
\end{frame}

%------------------------------------------------------------------------------%
\begin{proofframe}{O Subespaço Gerado Contém o Vetor Nulo}
  Seja
  \[
    W
    = \ger(G)
    \pause
    = \ger \left\{ \vec{v}_{1}, \ldots, \vec{v}_{k} \right\}
    \pause
    \subset V
  \]
  \pause
  O vetor nulo pertence a $W$
  \[
    \vec{0} \in W
  \]
  \pause
  pois
  \[
    \vec{0} = 0\vec{v}_{1} + \cdots + 0\vec{v}_{k}
  \]
\end{proofframe}

%------------------------------------------------------------------------------%
\begin{proofframe}{O Subespaço Gerado é Fechado}
  Sejam \(\vec{r}\) e \(\vec{s}\) vetores em
  \[
    W = \ger(G) = \ger \left\{ \vec{v}_{1}, \ldots, \vec{v}_{k} \right\}
  \]
  \pause
  precisamos mostrar que \(a\vec{r} + b\vec{s} \in W\)

  \pause
  Sabemos que
  \[
    \pause
    \vec{r} = c_1\vec{v}_1 + c_2\vec{v}_2 + \cdots + c_k\vec{v}_k
  \]
  \[
    \pause
    \vec{s} = d_1\vec{v}_1 + d_2\vec{v}_2 + \cdots + d_k\vec{v}_k
  \]
\end{proofframe}

%------------------------------------------------------------------------------%
\begin{proofframe}{O Subespaço Gerado é Fechado}
  \begin{align*}
    \onslide<+->{a\vec{r} + b\vec{s}}
    \onslide<+->{& = a\left(c_1\vec{v}_1 + c_2\vec{v}_2 + \cdots + c_k\vec{v}_k\right)}
    \onslide<+->{  + b\left(d_1\vec{v}_1 + d_2\vec{v}_2 + \cdots + d_k\vec{v}_k\right)}
    \\
    \onslide<+->{& = ac_1\vec{v}_1 + ac_2\vec{v}_2 + \cdots + ac_k\vec{v}_k}
    \onslide<+->{  + bd_1\vec{v}_1 + bd_2\vec{v}_2 + \cdots + bd_k\vec{v}_k}
    \\
    \onslide<+->{& = (ac_1+ bd_1)\vec{v}_1 + (ac_2 + bd_2)\vec{v}_2 + \cdots + (ac_k + bd_k)\vec{v}_k}
    \\
    \onslide<+->{& = p_1\vec{v}_1 + p_2\vec{v}_2 + \cdots + p_k\vec{v}_k}
  \end{align*}
  \onslide<+->{Então \(a\vec{r} + b\vec{s} \in W\)}
\end{proofframe}

%------------------------------------------------------------------------------%
\section{Solução de Sistemas Lineares}

%------------------------------------------------------------------------------%
\begin{frame}{Espaço Solução}
  Considere um sistema homogêneo
  \[
    A\vec{x} = \vec{0}
  \]
  \pause
  O conjunto solução é
  \[
    \pause
    W = \left\{ \vec{x}\in\R^n \st A\vec{x} = \vec{0} \right\}
  \]
  \pause
  Esse conjunto é sempre um subespaço de $\R^n$
\end{frame}

%------------------------------------------------------------------------------%
\begin{proofframe}{O Conjunto Solução Contém o Vetor Nulo}
  Seja
  \[
    W = \left\{ \vec{x}\in\R^n \st A\vec{x} = \vec{0} \right\}
    \subset \R^n
  \]
  \pause
  O vetor nulo pertence a $W$
  \[
    \vec{0} \in W
  \]
  \pause
  pois
  \[
    A \vec{0} = \vec{0}
  \]
\end{proofframe}

%------------------------------------------------------------------------------%
\begin{proofframe}{O Conjunto Solução é Fechado}
  Sejam \(\vec{r}\) e \(\vec{s}\) vetores em
  \[
    W = \left\{ \vec{x}\in\R^n \st A\vec{x} = \vec{0} \right\}
  \]
  \pause
  precisamos mostrar que \(a\vec{r} + b\vec{s} \in W\)

  \pause
  Sabemos que
  \[
    \pause
    A \vec{r} = \vec{0}
    \qquad\qquad
    \pause
    A \vec{s} = \vec{0}
  \]
  \pause
  Portanto
  \[
    \pause
    A\left(a\vec{r} + b\vec{s}\right)
    \pause
    = aA\vec{r} + bA\vec{s}
    \pause
    = a\vec{0} + b\vec{0}
    \pause
    = \vec{0}
  \]
\end{proofframe}

%------------------------------------------------------------------------------%
%------------------------------------------------------------------------------%
\begin{question}
  Determine se
  \[
    W = \left\{ \vetor{x}{y} \in\R^{2} \st y = 2x \right\}
  \]
  é um subespaço de $\R^{2}$
\end{question}

%------------------------------------------------------------------------------%
\begin{resolution}{Vetor Nulo}
  O vetor nulo \(\vetor{0}{0}\) pertence a $W$?

  \pause
  Se $x=0$ então
  \[
    \pause
    y
    \pause
    = 2x
    \pause
    = 2\times 0
    \pause
    = 0
  \]

  \pause
  Portanto
  \[
    \vec{0} \in W
  \]

\end{resolution}
%------------------------------------------------------------------------------%
\begin{resolution}{Combinação Linear}
  A combinação linear de dois vetores em $W$ pertence a $W$?
  \[
    \pause
    \vec{v}, \; \vec{w} \; \in W
    \pause
    \quad\Rightarrow\quad
    a \vec{v} + b\vec{w} \; \in W
  \]

  \pause
  Se \(\vec{v}, \vec{w} \in W\)
  \pause
  então
  \(
    \vec{v} = \vetor{p}{2p}
    \quad
    \vec{w} = \vetor{q}{2q}
  \)
  \pause
  \[
    a\vec{v}+b\vec{w}
    \pause
    = a\vetor{p}{2p} + b\vetor{q}{2q}
    \pause
    = \vetor{ap+bq}{2(ap+bq)}
    \pause
    = \vetor{x}{2x}
  \]
   \pause
  Logo
  \(
    a\vec{u}+b\vec{v}\in W
  \)

  \pause
  Portanto, $W$ é subespaço de $\R^2$
\end{resolution}

%------------------------------------------------------------------------------%
\endinput

%------------------------------------------------------------------------------%
\begin{question}
  Determine se
  \[
    W = \left\{ \vetor{x}{y} \in\R^{2} \st y=2x+1 \right\}
  \]
  é um subespaço de $\R^{2}$
\end{question}

%------------------------------------------------------------------------------%
\begin{resolution}{Vetor Nulo}
  O vetor nulo \(\vetor{0}{0}\) pertence a $W$?

  \pause
  Se $x=0$ então
  \[
    \pause
    y
    \pause
    = 3x + 1
    \pause
    = 3 \times 0 + 1
    \pause
    = 1
    \pause
    \neq 0
  \]

  \pause
  Logo
  \(
    \vec{0} \notin W
  \)

  \pause
  Portanto, $W$ não é subespaço de $\R^{2}$
\end{resolution}

%------------------------------------------------------------------------------%
\endinput
%------------------------------------------------------------------------------%
\begin{question}
  Determine se
  \[
    W = \left\{ \vetortri{x}{y}{z} \in\R^{3} \st x+y-z=0 \right\}
  \]
  é um subespaço de $\R^{3}$
\end{question}

%------------------------------------------------------------------------------%
\begin{resolution}{Vetor Nulo}
  O vetor nulo \(\vetortri{0}{0}{0}\) pertence a $W$?

  \pause
  \[
    x + y - z
    \pause
    = 0 + 0 - 0
    \pause
    = 0
  \]

  \pause
  Portanto
  \[
    \vec{0} \in W
  \]
\end{resolution}

%------------------------------------------------------------------------------%
\begin{resolution}{Combinação Linear}
  A combinação linear de dois vetores em $W$ pertence a $W$?
  \[
    \pause
    \vec{v}, \; \vec{w} \; \in W
    \pause
    \quad\Rightarrow\quad
    a \vec{v} + b\vec{w} \; \in W
  \]

  \pause
  Se \(\vec{v} = \vetortri{v_1}{v_2}{v_3}\) e \(\vec{w} = \vetortri{w_1}{w_2}{w_3}\in W\)

  \pause
  \[
    \vec{p}
    = a\vec{v}+b\vec{w}
    \pause
    = a \vetortri{v_1}{v_2}{v_3} + b \vetortri{w_1}{w_2}{w_3}
    \pause
    = \vetortri{av_1+bw_1}{av_2+bw_2}{av_3+bw_3}
  \]
\end{resolution}

%------------------------------------------------------------------------------%
\begin{resolution}{Combinação Linear}
  \onslide<+->{Como \(\vec{v}, \vec{w} \in W\)}
  \[
    \onslide<+->{v_1 + v_2 - v_3 = 0}
    \qquad
    \onslide<+->{w_1 + w_2 - w_3 = 0}
  \]
  \onslide<+->{Verificando se \(\vec{p} \in W\)}
  \begin{align*}
    \onslide<+->{x + y - z}
    \onslide<+->{& = p_1 + p_2 - p_3                         \\}
    \onslide<+->{& = (av_1+bw_1) + (av_2+bw_2) - (av_3+bw_3) \\}
    \onslide<+->{& = a(v_1+v_2-v_3) + b(w_1+w_2-w_3)         \\}
    \onslide<+->{& = a\times 0 + b\times 0                   \\}
    \onslide<+->{& = 0}
  \end{align*}
\end{resolution}

%------------------------------------------------------------------------------%
\begin{resolution}{Combinação Linear}
  \onslide<+->{Logo
  \(
    \vec{p}\in W
  \)}

  \onslide<+->{Portanto $W$ é subespaço de $\R^3$}
\end{resolution}

%------------------------------------------------------------------------------%
\endinput

%------------------------------------------------------------------------------%
\begin{question}
  Determine se
  \[
    W = \left\{ \vetortri{x}{y}{z} \in\R^{3} \st x+y-z = 2 \right\}
  \]
  é um subespaço de $\R^{3}$
\end{question}

%------------------------------------------------------------------------------%
\begin{resolution}{Vetor Nulo}
  O vetor nulo \(\vetortri{0}{0}{0}\) pertence a $W$?

  \pause
  \[
    x + y - z
    \pause
    = 0 + 0 - 0
    \pause
    = 0
    \pause
    \neq 2
  \]

  \pause
  Portanto
  \[
    \vec{0} \notin W
  \]

  \pause
  Portanto, $W$ não é subespaço de $\R^{2}$
\end{resolution}

%------------------------------------------------------------------------------%
\endinput

%%------------------------------------------------------------------------------%
\begin{question}
  Determine uma descrição paramétrica do subespaço
  \[
  W=
  \left\{
  \begin{pmatrix}
  x\\y\\z
  \end{pmatrix}
  \in\R^{3}
  \;\middle|\;
  2x-y+z=0
  \right\}
  \]
\end{question}

%------------------------------------------------------------------------------%
\begin{resolution}{Solução}
  Da equação
  \begin{align*}
  2x-y+z&=0\\
  z&=-2x+y
  \end{align*}
  Tomando
  \[
  x=s
  \qquad
  y=t
  \]
  obtemos
  \begin{align*}
  \begin{pmatrix}
  x\\y\\z
  \end{pmatrix}
  &=
  \begin{pmatrix}
  s\\t\\-2s+t
  \end{pmatrix}\\
  &=
  s
  \begin{pmatrix}
  1\\0\\-2
  \end{pmatrix}
  +
  t
  \begin{pmatrix}
  0\\1\\1
  \end{pmatrix}
  \end{align*}
  Logo
  \[
  W=
  \operatorname{span}
  \left\{
  \begin{pmatrix}
  1\\0\\-2
  \end{pmatrix}
  ,
  \begin{pmatrix}
  0\\1\\1
  \end{pmatrix}
  \right\}
  \]

\end{resolution}

%------------------------------------------------------------------------------%
\endinput

%------------------------------------------------------------------------------%
\begin{question}
  Determine o espaço solução de
  \[
  \begin{cases}
    x + y + z = 0 \\
    x - y + z = 0
  \end{cases}
  \]
\end{question}

%------------------------------------------------------------------------------%
\begin{resolution}{Solução}
\begin{columns}
\column{0.5\textwidth}
  \onslide<+->{Subtraindo as equações}
  \begin{align*}
    \onslide<+->{(x+y+z)-(x-y+z) & = 0 \\}
    \onslide<+->{2y              & = 0 \\}
    \onslide<+->{y               & = 0}
  \end{align*}
  \onslide<+->{Substituindo}
  \begin{align*}
    \onslide<+->{x + y + z & =  0  \\}
    \onslide<+->{x + z     & =  0  \\}
    \onslide<+->{x         & = -z}
  \end{align*}
  \column{0.5\textwidth}
  \onslide<+->{Tomando $z=t$}
  \[
    \onslide<+->{\vetortri{x}{y}{z}}
    \onslide<+->{= \vetortri{-t}{0}{t}}
    \onslide<+->{= t \vetortri{-1}{0}{1}}
  \]
  \onslide<+->{Logo
  \[
    W = \ger \left\{ \vetortri{-1}{0}{1} \right\}
  \]}
\end{columns}
\end{resolution}

%------------------------------------------------------------------------------%
\endinput

%------------------------------------------------------------------------------%
\begin{question}
  Determine se o conjunto das matrizes simétricas de ordem 2
  \[
    W=
    \left\{
    \begin{pmatrix}[cc]
      a & b \\
      b & c
    \end{pmatrix}
    \st
    a, b, c \in \R
    \right\}
  \]
  é um subespaço de $M_{2\times2}$
\end{question}

%------------------------------------------------------------------------------%
\begin{resolution}{Matriz Nula}
  Escolhendo $a$, $b$ e $c$ iguais a zero

  \pause
  verificamos que a matriz nula pertence a $W$
  \[
    \pause
    \begin{pmatrix}
      0 & 0 \\
      0 & 0
    \end{pmatrix}
    \in W
  \]
\end{resolution}

%------------------------------------------------------------------------------%
\begin{resolution}{Combinação Linear}
  Sejam
  \[
    A =
    \begin{pmatrix}
      a & b \\
      b & c
    \end{pmatrix}
    \qquad
    \pause
    B =
    \begin{pmatrix}
      d & e \\
      e & f
    \end{pmatrix}
    \qquad
    \pause
    \lambda, \mu\in\R
  \]
  \pause
  \[
    \lambda A + \mu B
    \pause
    = \lambda
    \begin{pmatrix}
      a & b \\
      b & c
    \end{pmatrix}
    + \mu
    \begin{pmatrix}
      d & e \\
      e & f
    \end{pmatrix}
    \pause
    =
    \begin{pmatrix}
      \lambda a+\mu d & \lambda b+\mu e \\
      \lambda b+\mu e & \lambda c+\mu f
    \end{pmatrix}
  \]

  \pause
  A matriz resultante é simétrica

  \pause
  Logo $W$ é subespaço
\end{resolution}

%------------------------------------------------------------------------------%
\endinput

%%------------------------------------------------------------------------------%
\begin{question}
  Determine se
  \[
  W=
  \left\{
  f:\R\to\R
  \;\middle|\;
  f(0)=0
  \right\}
  \]
  é um subespaço do espaço de funções reais
\end{question}

%------------------------------------------------------------------------------%
\begin{resolution}{Solução}
  A função nula satisfaz
  \begin{align*}
  0(0)&=0
  \end{align*}
  Logo
  \[
  0\in W
  \]
  Sejam $f,g\in W$.
  Então
  \[
  f(0)=0
  \qquad
  g(0)=0
  \]
  Para $\lambda,\mu\in\R$
  \begin{align*}
  (\lambda f+\mu g)(0)
  &=\lambda f(0)+\mu g(0)\\
  &=\lambda\cdot0+\mu\cdot0\\
  &=0
  \end{align*}
  Logo
  \[
  \lambda f+\mu g\in W
  \]
  Portanto $W$ é subespaço.

\end{resolution}

%------------------------------------------------------------------------------%
\endinput

%%------------------------------------------------------------------------------%
\begin{question}
  Determine se
  \[
  W=
  \left\{
  \begin{pmatrix}
  x\\y\\z
  \end{pmatrix}
  \in\R^{3}
  \;\middle|\;
  xy=0
  \right\}
  \]
  é um subespaço de $\R^{3}$
\end{question}

%------------------------------------------------------------------------------%
\begin{resolution}{Solução}
  Temos
  \[
  \vec{u}=
  \begin{pmatrix}
  1\\0\\0
  \end{pmatrix}
  \in W
  \qquad
  \vec{v}=
  \begin{pmatrix}
  0\\1\\0
  \end{pmatrix}
  \in W
  \]
  Mas
  \begin{align*}
  \vec{u}+\vec{v}
  &=
  \begin{pmatrix}
  1\\1\\0
  \end{pmatrix}\\
  1\cdot1&=1\neq0
  \end{align*}
  Logo
  \[
  \vec{u}+\vec{v}\notin W
  \]
  Portanto $W$ não é subespaço
\end{resolution}

%------------------------------------------------------------------------------%
\endinput


%------------------------------------------------------------------------------%
%\section{Lista Mínima}
%
%%------------------------------------------------------------------------------%
%\begin{frame}{Lista Mínima}
%  \begin{enumerate}
%    \item
%  \end{enumerate}
%
%  \vfill
%  Atenção: A prova é baseada no livro, não nas apresentações
%\end{frame}

%------------------------------------------------------------------------------%
\end{document}
%------------------------------------------------------------------------------%
,